\documentclass{article}
\usepackage{graphicx} % Required for inserting images
%\usepackage[a4paper, total={6in, 8in}]{geometry}
\usepackage{amsmath, amssymb, empheq, mathtools, bbold, enumitem, hyperref, physics}
\usepackage{amsthm}
\usepackage[framemethod=TikZ]{mdframed}
%\usepackage{framed}
%\usepackage[amsthm,framed,thmmarks]{ntheorem}

\renewcommand{\bf}[1]{\textbf{#1}}
\renewcommand{\v}[1]{\boldsymbol{#1}}

%\newframedtheorem{frm-thm}{Theorem}
\theoremstyle{definition}
\newtheorem{theorem}{Theorem}
\newtheorem*{fact}{Fact}
\newtheorem*{note}{Note}
\newtheorem*{remark}{Remark}
\newtheorem{definition}{Definition}

\surroundwithmdframed[]{theorem}
\newenvironment{frmthm}{
	\begin{mdframed}
		\begin{theorem}
		}{		
		\end{theorem}
	\end{mdframed}
}

\newenvironment{factsummary}{
	{\large\textbf{Fact Summary}}\\
	}{
}
\surroundwithmdframed{factsummary}

\newcommand{\xlrightarrow}{\xrightarrow{\hspace{0.8cm}}}
\newcommand{\reals}{\mathbb{R}}

\newcommand{\matadd}{+_{\text{mat}}}
\newcommand{\matdot}{\cdot_{\text{mat}}}
\newcommand{\vp}{\oplus}
\newcommand{\vm}{\odot}

\newcommand{\mxn}{m \times n}
\newcommand{\letmatrix}[2]{\mathcal{M}_{#1 \times #2}(\reals)}
\newcommand{\rowop}{\mathcal{R}}
%\newcommand{\span}[1]{\text{span(}#1\text{)}}


\newcommand{\bmat}[1]{\begin{bmatrix}#1\end{bmatrix}}
\newcommand{\detmat}[1]{
	\left|\hspace{0.5em}
	\begin{matrix}
		#1
	\end{matrix}
	\hspace{0.5em}\right|
}

% Ch. 3: Vector Spaces

\begin{document}
\section{Definition of a Vector Space}
	\begin{definition}[Vector Space]
		A set $V$ is a \underline{vector space over $\reals$} if there is an operation "$\oplus$" combining $\alpha, \beta \in V$ into $\alpha \oplus \beta \in V$, and an operation "$\odot$" combining $c \in \reals$ and $\alpha \in V$ into $c \odot \alpha \in V$, which furthermore satisfy the following axioms for all $\alpha, \beta, \delta \in V$ and $c,d \in \reals$:
		
		\begin{enumerate}[font=\bfseries, itemsep=0pt]
			\item The sum $\alpha \oplus \beta$ is in $V$. \hfill Closed under addition
			\item $\alpha \oplus \beta = \beta \oplus \alpha$ \hfill Addition is commutative
			\item $(\alpha \oplus \beta) \oplus \delta = \alpha \oplus (\beta \oplus \delta)$ \hfill Addition is associative
			\item There exists a vector $\vb{0}_v \in V$ s.t. for every vector $\alpha \in V, \vb{0}_v \oplus \alpha = \alpha$. \hfill Additive identity
			\item For every vector $\alpha \in V$, there exists a vector, denoted by $\delta$, such that $\alpha \oplus \delta = \vb{0}_v$. \hfill Additive inverse
			\item The scalar product $c \odot \alpha$ is in $V$. \hfill Closed under scalar multiplication
			\item $c \odot (\alpha \oplus \beta) = (c \odot \alpha) \oplus (c \odot \beta)$
			\item $(c + d) \odot \alpha = (c \odot \alpha) \oplus (d \odot \alpha)$
			\item $c \odot (d \odot \alpha) = (cd)\odot \alpha$
			\item $1 \odot \alpha = \alpha$
		\end{enumerate}
	\end{definition}

	\begin{remark}
		$\text{dim}(V) = n \implies V \cong \reals^n$
	\end{remark}

	\begin{factsummary}
		\begin{enumerate}[font=\bfseries,itemsep=0pt]
			\item To determine whether a set $V$ with addition and scalar multiplication defined on $V$ is a vector space requires verification of the 10 vector space axioms.
			\item The Euclidean space $\reals^n$ and the set of matrices $M_{m\times n}$, with the standard componentwise operations, are vector spaces. The set of polynomials of degree $n$ or less, including the zero polynomial, with termwise operations is a vector space.
			\item In all vector spaces, additive inverses are unique. Also
			\[ 0 \odot \alpha = \vb{0}_v \qquad c \cdot \vb{0}_v = \vb{0}_v \qquad \text{and} \qquad (-1)\odot \alpha = -\alpha \]
			In addition if $c \odot \vb{u} = \vb{0}$, then either the scalar $c$ is the number 0 or the vector $\vb{u}$ is the zero vector.
		\end{enumerate}
	\end{factsummary}
\section{Subspaces}
	\begin{definition}[Subspace]
			A \textbf{subspace} $W$ of a vector space $V$ is a nonempty subset that is itself a vector space with respect to the inherited operations of vector addition and scalar multiplication on $V$.
	\end{definition}
	\begin{theorem}
		Let $S$ be a nonempty subset of the vector space $V$. Then $S$ is a subspace of $V$ $\iff$ $S$ is closed under addition and scalar multiplication.
	\end{theorem}
	\begin{theorem}
		A nonempty subset $S$ of a vector space $V$ is a subspace of $V$ $\iff$ each pair of vectors $\alpha$ and $\beta$ in $S$ and each scalar $c$, the vector $\alpha \oplus (c \odot \beta)$.
	\end{theorem}
	\begin{theorem}
		The intersection of any collection of subspaces of a vector space is a subspace of the vector space.
	\end{theorem}
	\subsection{Span of a Set of Vectors}
	\begin{definition}[Linear Combination]
		Let $S = \{\alpha_1, \alpha_2, ..., \alpha_k\}$ be a set of vectors in a vector space $V$, and let $c_1, c_2, ..., c_k$ be scalars. A \textbf{linear combination} of the vectors of $S$ is an expression of the form
		\[ (c_1 \odot \alpha_1) \oplus (c_2 \odot \alpha_2) \oplus \cdots \oplus (c_k \odot \alpha_k) \]
	\end{definition}
	\begin{definition}[Span of a Set of Vectors]
		Let $V$ be a vector space and let $S = \{\alpha_1, \alpha_2, ..., \alpha_k\}$
		be a set of vectors in $V$. The \textbf{span} of $S$, denoted by $\text{span}(S)$, is the set
		\[ \text{span}(S)= \{ c_1\alpha_1 = c_2\alpha_2 + \cdots + c_n\alpha_n | c_1, c_2, ..., c_n \in \reals \} \]
		$S$ is a \emph{spanning set for $V$} if span($S$) = $V$.
	\end{definition}
	\begin{theorem}
		If $S \{\alpha_1, \alpha_2, ..., \alpha_n\}$ is a set of vectors in a vector space $V$, then \bf{span}($S$) is a subspace.
	\end{theorem}
	\begin{definition}[Null Space and Column Space]
		Let $A$ be an $m \times n$ matrix.
		\begin{itemize}[font=\bfseries,itemsep=0pt]
			\item The \bf{null space} of $A$, denoted by $N(A)$, is the set of all vectors in $\reals^n$ such that $A\bf{x} = \bf{0}$.
			\item The \bf{column space} of $A$, denoted by $col(A)$, is the set of all linear combinations of the column vectors of $A$.
		\end{itemize}
	\end{definition}
	\begin{theorem}
		Let $A$ be an $m \times n$ matrix. The linear system $A\bf{x} = \bf{b}$ is consistent $\iff$ \bf{b} is in the column space of $A$.
	\end{theorem}
	\begin{theorem}
		Let $A$ be an $m \times n$ matrix. Then the null space of $A$ is a subspace of $\reals^n$.
	\end{theorem}
\section{Basis and Definition}
	\begin{definition}[Basis for a Vector Space]
		A subset $B$ of a vector space $V$ is a \textbf{basis} for $V$ provided that
		\begin{enumerate}[font=\bfseries,itemsep=0pt]
			\item $B$ is a linearly independent set of vectors in $V$
			\item span($B$) = $V$
		\end{enumerate}
		In other words, a \bf{basis} for a v.sp. $V$ is a linearly independent subset $S \subseteq V$ for which $span(S) = V$.
	\end{definition}
	\begin{theorem}
		If $\{\alpha_1, ..., \alpha_k\}$ is lin. ind., and if $c_1\alpha_1 + \cdots + c_k\alpha_k = d_1\alpha_1 + \cdots + d_k\alpha_k$, then $c_i = d_i$ for all $i$.
	\end{theorem}
	Let $\gamma = \{(1,1,1),(1,1,0),(1,2,3)\} = \{f_1,f_2,f_3\}$ as an ordered basis for $\reals^3$. and $\alpha = (3,4,5) = c_1f_1 + c_2f_2 + c_3f_3 = c_1(1,1,1) + c_2(1,1,0) + c_3(1,2,3)$.\\
	\[ \bmat{1&1&1&3\\1&1&2&4\\1&0&3&5} \longrightarrow \bmat{1&1&1&3\\0&0&1&1\\0&-1&2&2} \longrightarrow \bmat{1&1&1&3\\0&1&-2&-2\\0&0&1&1}\]
	Use back substitution to find solutions:
	\begin{align*}
		c_3 = 1\\
		c_2 = -2 + 2(1) = 0\\
		c_1 = 3 - 1(0) - 1(1) = 2
	\end{align*}
	\[ [\alpha]_\gamma = \bmat{2\\0\\1} \text{, in others words } \bmat{3\\4\\5}_\gamma = \bmat{2\\0\\1} \]
	(said aloud as \emph{alpha with respect to gamma}).
\end{document}